\uniterablesection{Вступ}

Формування та утримання корисних звичок — одна з визначальних умов особистої продуктивності й професійного зростання. Особливо помітно це у сфері розробки програмного забезпечення, де щоденна практика — написання коду, рецензування чужих напрацювань, читання документації, фізична активність для збереження працездатності — безпосередньо позначається на кваліфікації фахівця. Дослідники, які вивчають механізми поведінки людини, наголошують, що сталий результат дають не разові зусилля, а саме регулярні, нехай і невеликі, повторювані дії \citecustom{1}, \citecustom{2}.

На ринку представлено чимало мобільних і вебзастосунків для трекінгу звичок, проте переважна їх частина побудована на одному й тому самому принципі — ручному відмічанні виконаних дій. Такий підхід має суттєву ваду: він вимагає від людини дисциплінованості саме в момент фіксації результату. Через це зібрані дані нерідко виявляються неповними або суб'єктивними, а сам облік перетворюється на ще одну звичку, яку так само важко підтримувати. Виникає парадоксальна ситуація, коли інструмент, покликаний полегшити роботу над собою, додає користувачеві навантаження.

Водночас значна частка активності сучасного розробника вже фіксується зовнішніми сервісами в автоматичному режимі. Платформа GitHub зберігає історію комітів та запитів на злиття, сервіс Strava — тренування й іншу фізичну активність. Якщо під'єднати ці джерела до системи трекінгу звичок, відмічання прогресу можна зробити автоматичним: коміт у репозиторій чи завершене тренування самі стають свідченням виконаної звички. Це знижує навантаження на користувача й одночасно підвищує достовірність зібраних даних. Саме на подоланні розриву між реальною активністю людини та її обліком і зосереджено цю роботу.

\textbf{Мета роботи} — спроєктувати й розробити інформаційну систему моніторингу особистих звичок, яка поєднує ручне ведення записів з автоматичним збором даних із зовнішніх сервісів (GitHub, Strava) та надає користувачеві наочні засоби аналізу власного прогресу.

Для досягнення поставленої мети потрібно розв'язати такі завдання:

\begin{enumerate}[noitemsep, topsep=2pt, leftmargin=*, label=\arabic*)]
    \item проаналізувати предметну область трекінгу звичок, дослідити наявні рішення-аналоги та виявити їхні обмеження;
    \item сформулювати функціональні й нефункціональні вимоги до системи;
    \item обґрунтувати вибір архітектурного підходу та технологічного стеку;
    \item спроєктувати доменну модель, структуру бази даних і програмний інтерфейс;
    \item реалізувати серверну частину з модульною архітектурою та механізмом автоматичного опитування зовнішніх сервісів;
    \item розробити клієнтський вебзастосунок з наочною візуалізацією активності;
    \item провести тестування ключових компонентів системи.
\end{enumerate}

\textbf{Об'єкт дослідження} — процес формування, обліку та аналізу особистих звичок користувача.

\textbf{Предмет дослідження} — методи та програмні засоби автоматизованого збору й обробки даних про активність користувача для моніторингу його звичок.

\textbf{Методи дослідження.} У роботі застосовано методи системного аналізу — для вивчення предметної області та декомпозиції задачі; принципи предметно-орієнтованого проєктування (DDD) і чистої архітектури — для побудови доменної моделі та поділу системи на шари; методи об'єктно-орієнтованого програмування й модульного тестування — для реалізації та перевірки коректності окремих компонентів.

\textbf{Практичне значення.} Розроблена система дає змогу розробникам та іншим користувачам об'єктивно відстежувати власні звички з мінімальними витратами зусиль на ручний облік. Закладені архітектурні рішення — модульний моноліт, розділення команд і запитів (CQRS), подієва взаємодія між модулями — залишають простір для подальшого розширення системи новими інтеграціями та її масштабування без істотного переписування коду.

\textbf{Структура роботи.} Робота складається зі вступу, трьох розділів, висновків, списку використаних джерел і додатків. У першому розділі проаналізовано предметну область і сформульовано вимоги до системи. Другий розділ присвячено проєктуванню архітектури, доменної моделі та бази даних, а також обґрунтуванню обраних технологій. У третьому розділі описано реалізацію серверної й клієнтської частин, механізм інтеграції із зовнішніми сервісами та тестування.
